\input{build/foundations-preamble.tex}
\usepackage[authoryear]{natbib}
\renewcommand{\refname}{સંદર્ભગ્રંથો}
\renewcommand{\partname}{ભાગ}
\newenvironment{intro}{}{}
\newcommand{\Id}[1]{\mathord{\mathrm{Id}_{#1}}}
\newcommand{\equivrep}[2]{[#1]_{#2}}
\newcommand{\equivclass}[2]{#1/_{\!{#2}}}
\newcommand{\Intequiv}{\mathrel{\sim_{\Int}}}
\newcommand{\Ratequiv}{\mathrel{\sim_{\Rat}}}
\newcommand{\Realequiv}{\mathrel{\sim_{\Real}}}
\newcommand{\funrestrictionto}[2]{#1\mathord{\restriction}_{#2}}
\newcommand{\funimage}[2]{#1[#2]}
\newcommand{\emptyseq}{\Lambda}
\newcommand{\liff}{\mathbin{\leftrightarrow}}
\newcommand{\dom}[1]{\mathrm{dom}(#1)}
\newcommand{\ran}[1]{\mathrm{ran}(#1)}
\newcommand{\comp}[2]{#2\circ #1}
\newcommand{\pto}{\mathrel{\ooalign{\hfil$\mapstochar\mkern5mu$\hfil\cr$\to$}}}
\newcommand{\fdefined}{\downarrow}
\newcommand{\fundefined}{\uparrow}
\newcommand{\sourcecorrection}[2]{%
  \par\smallskip\noindent
  \fbox{\begin{minipage}{0.94\linewidth}\small
    \textbf{સ્રોત-સુધારો #1.} #2
  \end{minipage}}%
  \par\smallskip}
\definecolor{oldiagcolorD}{HTML}{1973ba}
\newlength{\olphotowidth}
\setlength{\olphotowidth}{.4\textwidth}
\RenewDocumentCommand{\olasset}{O{\olphotowidth} m}{\centering\resizebox{#1}{!}{\input{upstream/#2}}}
\input{build/arithmetization-available-labels.tex}
\renewcommand{\oliflabeldef}[3]{\ifcsname guavailable@#1\endcsname#2\else#3\fi}
\hypersetup{pdftitle={ગણો, સંબંધો, વિધેયો, ગણોનું કદ અને અંકગણિતીકરણ — ઓપન લોજિક ગુજરાતી}}
\begin{document}
\begin{titlepage}
{\Huge\bfseries ગણો, સંબંધો, વિધેયો, ગણોનું કદ અને અંકગણિતીકરણ\par}
\vspace{1em}
{\Large ઓપન લોજિક — ગુજરાતી\par}
\vspace{2em}
ઓપન લોજિક પ્રોજેક્ટના અંગ્રેજી મૂળનો ગુજરાતી અનુવાદ.\par
\vspace{1em}
આ આવૃત્તિમાં ગણો, સંબંધો, વિધેયો, ગણોનું કદ અને અંકગણિતીકરણનાં
સંપૂર્ણ પ્રકરણો છે. પ્રાકૃતિક સંખ્યાઓમાંથી પૂર્ણાંક અને સંમેય સંખ્યાઓ,
ડેડેકિન્ડ કાપ તથા કોશી શ્રેણીઓ દ્વારા વાસ્તવિક સંખ્યાઓની રચના,
પૂર્ણતા ગુણધર્મ અને સંબંધિત તાત્ત્વિક ચર્ચા સતત ક્રમમાં સમાવ્યાં છે.\par
\vfill
આ યંત્ર દ્વારા કરેલો અનુવાદ છે. સંપૂર્ણ ગ્રંથનું કામ ચાલુ છે.
આ આવૃત્તિમાં મૂળનાં ૪૫ એકમોનો સમાવેશ છે; કુલ ૭૨૨ એકમો છે.
કેટલીક વિશિષ્ટ પરિભાષા કામચલાઉ છે. ચકાસણીની વિગતો અને
મૂળ પાઠ વિશેની નોંધો સંપાદકીય માહિતીમાં આપેલી છે.\par
\vspace{1em}
\href{https://github.com/OpenLogicProject/OpenLogic}{Open Logic Project}
\quad \href{https://creativecommons.org/licenses/by/4.0/}{CC BY 4.0}\par
\href{https://github.com/KokunoYumeto/OpenLogic-translations}{અનુવાદોનું કેન્દ્ર}\par
\end{titlepage}
\tableofcontents
\clearpage
\input{build/arithmetization-body.tex}
\clearpage
\section*{સંપાદકીય નોંધો}
\addcontentsline{toc}{section}{સંપાદકીય નોંધો}
આ નોંધો મૂળ પાઠથી અલગ છે. મૂળની ગણિતીય નિશાનીઓ જાળવી છે.
OLFUN-001થી OLFUN-005, OLSIZ-001થી OLSIZ-012 અને OLARI-001થી
OLARI-004 સુધીની ઓળખવાળી નોંધો સ્થિર અંગ્રેજી મૂળમાં મળેલી ખામીઓ અને
ગુજરાતી પાઠમાં કરેલા મર્યાદિત સુધારા બાજુબાજુ દર્શાવે છે. દરેક
સુધારાની ચોક્કસ જગ્યા, કારણ અને સૂત્રમાં થયેલો ફેરફાર સાથેના
નિર્ણય-અભિલેખમાં નોંધાયેલ છે.

અંકગણિતીકરણમાં પૂર્ણાંક અને સંમેય સંખ્યાઓ સામ્ય વર્ગો દ્વારા અને
વાસ્તવિક સંખ્યાઓ ડેડેકિન્ડ કાપો તથા વૈકલ્પિક રીતે કોશી શ્રેણીઓના
સામ્ય વર્ગો દ્વારા રચાય છે. આ ગણસૈદ્ધાંતિક રજૂઆતો રોજિંદા ગણિતીય
વ્યવહારને બદલતી ઓળખો તરીકે નહીં, પણ સંબંધિત સંરચનાઓનાં પ્રતિરૂપ તરીકે
વાંચવાની છે. મૂળમાં સામ્ય સંબંધ, સામ્ય વર્ગ અને સંખ્યાસંહતિના પ્રકાર
વચ્ચે થયેલા ચાર સ્પષ્ટ ભેદભંગો અનુવાદમાં નોંધ સાથે સુધાર્યા છે.

ગણનીય માટેનું ગુજરાતી પ્રમાણ ગણનીયને ગણનીય અનંત ગણ માટે વાપરે છે,
જ્યારે Open Logicની આ આવૃત્તિની વ્યાખ્યા સાન્ત અને ખાલી ગણોને પણ
સમાવે છે. બંને પરિગણના-રૂપો પોતાની સ્થાનિક વ્યાખ્યા સાથે જ વાંચવાનાં છે.
પ્રાથમિક વિભાગો અને વધુ અમૂર્ત વૈકલ્પિક વિભાગો બંને સાથે મૂક્યાં છે,
અને સશરતી પાઠ ઉપલબ્ધ વિભાગોને અનુરૂપ રીતે પસંદ થાય છે.

પરિભાષા, મૂળ સ્રોતો, પ્રત્યેક અનુવાદિત ખંડની સ્રોત-સંગતિ અને
ચકાસણીની વધુ વિગતો આવૃત્તિની સાથેની ફાઇલોમાં છે.
\begin{thebibliography}{7}
\bibitem[Benacerraf(1965)]{Benacerraf1965}
Benacerraf, Paul. 1965. \emph{What numbers could not be}.
The Philosophical Review 74(1), 47--73.
\bibitem[Cantor(1892)]{Cantor1892}
Cantor, Georg. 1892. Über eine elementare Frage der Mannigfaltigkeitslehre.
\bibitem[Frege(1884)]{Frege1884}
Frege, Gottlob. 1884. \emph{Die Grundlagen der Arithmetik}.
\bibitem[Potter(2004)]{Potter2004}
Potter, Michael. 2004. \emph{Set Theory and Its Philosophy}.
\bibitem[Conway(2006)]{Conway2006}
Conway, John. 2006. \emph{The Power of Mathematics}.
\bibitem[Katz and Katz(2012)]{KatzKatz2012}
Katz, Karin Usadi and Mikhail G. Katz. 2012. Stevin Numbers and Reality.
\bibitem[O'Connor and Robertson(2005)]{OConnorRobertson:RN}
O'Connor, John J. and Edmund F. Robertson. 2005. The real numbers: Stevin to Hilbert.
\end{thebibliography}
\end{document}
